\PassOptionsToPackage{numbers,square,sort&compress}{natbib}
\documentclass[11pt,twocolumn]{article}

\IfFileExists{acl.sty}{
  \usepackage[review]{acl}
}{
  \usepackage[margin=0.72in,columnsep=0.24in]{geometry}
  \usepackage[numbers,square,sort&compress]{natbib}
}

\usepackage{times}
\usepackage{latexsym}
\usepackage[T1]{fontenc}
\usepackage[utf8]{inputenc}
\usepackage{microtype}
\usepackage{inconsolata}
\usepackage{graphicx}
\usepackage{booktabs}
\usepackage{amsmath}
\usepackage{float}
\usepackage{placeins}
\usepackage{capt-of}
\IfFileExists{xurl.sty}{\usepackage{xurl}}{\usepackage{url}}
\setcitestyle{numbers,square}
\makeatletter
\@ifpackageloaded{hyperref}{}{\usepackage[hidelinks]{hyperref}}
\makeatother

\graphicspath{{figures/}}

\title{Predicting Ideology from Social Media Content Using Weakly Supervised Transformer Models}

\author{Jacob Kim and Billy Fox \\
  Rochester Institute of Technology \\
  \texttt{jck9187@rit.edu, bcf7023@rit.edu}}
\date{}

\begin{document}
\maketitle

\begin{abstract}
This paper tests whether social media posts can be used to predict political ideology without reducing each post to one left-right label. The system is designed around several connected decisions: whether a post is politically relevant, what issue or actor it targets, what stance and frame it uses, and whether the ideological signal is economic, social, intense, or ambiguous. The completed real-data experiment uses MITweet with a balanced TF-IDF baseline and a pretrained RoBERTa baseline. Both models beat random performance, but the transformer result is uneven. RoBERTa improves social-direction prediction, while economic-direction prediction collapses toward the majority neutral class. The balanced TF-IDF baseline does better on economic direction because it recovers more minority-class signal. The later B5 and B6 experiments should be read as pipeline checks on a small smoke run, not as final evidence that the hierarchical model is better.
\end{abstract}

\section{Introduction}

Ideology detection is the task of predicting political viewpoints from text. This is useful because political discussion now happens at large scale on social media, where researchers may want to study public opinion, ideological polarization, and how communities react to policy debates. A simple classifier might ask whether a post is liberal or conservative. This project studies a harder version of that task. A post can be relevant or irrelevant to politics, aimed at a particular policy target, supportive or opposed to that target, framed around economic or social concerns, and still be mixed or ambiguous. In this paper, ``harder'' means that the model must predict several connected labels rather than only one overall ideology class.

The main research question is whether weakly supervised transformer models can predict multidimensional ideology from social media text while still working when topics, communities, or platforms change. MITweet motivates the question because it treats ideology as a set of facets rather than one axis \cite{liu-etal-2023-ideology}. Work on biased supervision gives the practical risk: a classifier can learn topic shortcuts instead of ideological reasoning when labels are limited or uneven \cite{chen-etal-2023-ideology}. In that failure mode, the model learns that a topic often appears in one political community, not that the post is taking a particular stance.

The planned system combines high-quality labels with larger weakly labeled data. MITweet provides carefully annotated labels for ideological relevance and ideological direction across facets. Reddit Political Compass provides weaker signals along economic and social axes. Reddit Politosphere provides broad political subreddit text and community structure, but it does not directly provide post-level ideology labels in the same form as MITweet. In our design, Politosphere is used as a source of political text and community shift, while weak labels are derived from community metadata, lexicons, and Political Compass supervision rather than treated as gold ideology labels.

MITweet is the completed real-data experiment because Politosphere access was blocked in the available environment. The Reddit branch therefore ran only as an infrastructure smoke check, not as a final Reddit-scale experiment. On MITweet, RoBERTa helps most when the label distribution gives it usable signal, which is the case for social direction. It struggles when one class dominates, which happens for economic direction. The smoke run also shows that the hierarchical B6 model is not automatically stronger than the flatter B5 model under limited training. The strongest conclusion is therefore narrower than the original plan: the pipeline works, the MITweet evidence is real, and the next full run needs imbalance-aware losses before the model ladder can be judged fairly.

\section{Related Work}

This project depends most directly on MITweet, Reddit Politosphere, Reddit Political Compass, and work on biased supervision. MITweet is the main gold-label source because it annotates ideology as a set of facets rather than a single party or left-right label \cite{liu-etal-2023-ideology}. That design shapes our prediction heads for relevance, economic direction, and social direction. It also gives a concrete way to test whether a model can learn ideological direction from short social media posts.

Reddit Politosphere is useful for a different reason. It is a large text and network resource covering political subreddit discussion over time \cite{hofmann-etal-2022-politosphere}. We do not treat Politosphere as if it already contains the same post-level ideology labels as MITweet. Instead, it supports the planned weak-supervision and robustness setting: train on broad Reddit political language, then test whether the model generalizes across communities. Reddit Political Compass is closer to our label space because it studies economic and social ideology axes on Reddit \cite{colacrai-etal-2024-politicalcompass}. That makes it a natural weak-supervision source for the economic-direction and social-direction heads.

Prior modeling work explains why these datasets should not be used only with random train-test splits. Chen et al. show that ideology prediction can overfit the topic of a post rather than the way ideology is expressed \cite{chen-etal-2023-ideology}. This directly motivates our planned leave-topic-out, leave-community-out, and cross-platform checks. RoBERTa is used as the transformer encoder because it is a strong general text representation model \cite{liu2019roberta}, but the results here show that a transformer still needs loss functions and sampling strategies that match the label distribution. Simpler lexical and sequence baselines also remain important because older political text classification work shows that non-transformer models can be competitive when lexical cues are strong \cite{chang-masterson-2020-wordorder}.

Disagreement-aware learning is also relevant. Work on learning with disagreements treats annotator variation as signal rather than noise \cite{leonardelli-etal-2025-lewidi}. That matters for political text because a post can be genuinely mixed, sarcastic, or unclear. In our setting, a hard label is the single majority class used for ordinary classification. A soft label is the full annotator distribution over classes. Soft labels let the model learn that a post split between support and mixed is different from a post where every annotator agreed.

\section{Methodology}

The model uses eight prediction heads because ideology in social media text is not only a final class label. Relevance identifies whether the post is political. Target identifies the issue or actor being discussed. Stance marks support, opposition, or a mixed position toward that target. Frame captures the kind of reasoning being used, such as economic, social, security, rights, or other framing. Economic direction and social direction place the post on two ideology axes. Intensity captures how strong the ideological signal is, while ambiguity marks posts where the label is uncertain, mixed, or likely to split annotators.

For the real-data experiment, MITweet labels were projected into the heads that the dataset can support cleanly. The local MITweet file contains 12,594 labeled posts. Relevance is positive when any MITweet ideology facet is marked relevant. Economic direction is derived from MITweet's economic facets, and social direction is derived from its social facets. If no facet on an axis is active, the projected label is neutral. This projection lets MITweet evaluate relevance, economic direction, and social direction, but it does not fully evaluate target, stance, frame, intensity, or ambiguity.

For the Reddit branch, the codebase implements ingestion, weak-label construction, annotation folding, counterfactual editing, model training, prediction export, calibration, paired bootstrap testing, and error analysis. Weak labels are derived from subreddit and user priors, keyword dictionaries for issue frames, and Reddit Political Compass economic and social axes. This is the part of the project that operationalizes combining datasets: MITweet supplies clean labels, Reddit Political Compass supplies weak economic and social signals, and Politosphere supplies broader political communities for robustness tests.

B1 through B6 are shorthand for the project model ladder. B1 is a TF-IDF logistic-regression baseline with balanced class weights. B2 is a RoBERTa-base classifier fine-tuned on the available MITweet labels with ordinary cross-entropy. B3 through B5 are flat transformer variants intended for weak pretraining and gold fine-tuning. B5 adds decomposition heads for target, stance, frame, intensity, and ambiguity. B6 is the hierarchical representation model. It maps text into a latent ideology representation and jointly predicts discourse and ideology heads. The purpose of B6 is to test whether modeling discourse structure helps ideology prediction, not just whether a larger model improves the score.

\section{Experimental Design}

The original experimental design evaluates whether the final model generalizes across topics, communities, and platforms. The completed real-data part compares B1 and B2 on MITweet with three random seeds. Macro-F1 is the main metric because the ideology labels are imbalanced and a model can look accurate while mostly predicting the dominant class. The comparison also uses paired bootstrap testing over shared validation posts, which checks whether one model still wins after repeatedly resampling the same validation set. Calibration is evaluated by comparing model confidence on correct and incorrect predictions.

Table~\ref{tab:mitweet} shows the real MITweet results. Both B1 and B2 beat random performance. B2 does slightly better on relevance and social direction, but B1 does better on economic direction. The economic-direction split explains why. About 91\% of the economic labels are neutral. Across seeds, B2 reaches high F1 for the neutral class but near-zero F1 for the left and right classes. B1 handles those minority classes better because it uses balanced class weights.

\begin{center}
\centering
\small
\begin{tabular}{lccc}
\toprule
Label & B1 TF-IDF & B2 RoBERTa & Random \\
\midrule
Relevance & $0.708 \pm 0.077$ & $0.743 \pm 0.064$ & $\sim 0.41$ \\
Economic & $0.459 \pm 0.021$ & $0.341 \pm 0.017$ & $\sim 0.21$ \\
Social & $0.534 \pm 0.015$ & $0.579 \pm 0.022$ & $\sim 0.34$ \\
\midrule
Mean & $0.567 \pm 0.033$ & $0.554 \pm 0.024$ & $0.314$ \\
\bottomrule
\end{tabular}
\captionof{table}{Three-seed MITweet results. Values are macro-F1 mean and standard deviation.}
\label{tab:mitweet}
\end{center}

\begin{center}
    \centering
    \includegraphics[width=\columnwidth]{figure01_model_ladder_heatmap.pdf}
    \captionof{figure}{Heatmap of supported-label macro-F1 for the model ladder. The completed MITweet comparison is the reliable real-data result, while later ladder entries are diagnostic artifacts from the smoke-run branch.}
    \label{fig:ladder}
\end{center}

The paired bootstrap comparison gives the same interpretation. B1 significantly beats B2 on economic direction for all three seeds. On social direction, B2 beats B1 for one seed and ties for the other two. This means RoBERTa is useful when the label distribution gives it enough signal, but it is not robust to imbalance by default.

\begin{center}
\centering
\scriptsize
\begin{tabular}{lccc}
\toprule
Label & S42 & S43 & S44 \\
\midrule
Relevance & Tie (.47) & Tie (.55) & Tie (.36) \\
Economic & B1 $>$ B2 (.008) & B1 $>$ B2 (.014) & B1 $>$ B2 ($<$.001) \\
Social & Tie (.48) & B2 $>$ B1 ($<$.001) & Tie (.26) \\
\bottomrule
\end{tabular}
\captionof{table}{Paired bootstrap outcomes for B1 and B2 on MITweet validation posts.}
\label{tab:bootstrap}
\end{center}

\begin{center}
\centering
\small
\begin{tabular}{lccc}
\toprule
Label & Correct & Error & Gap \\
\midrule
Relevance & 0.957 & 0.849 & +0.107 \\
Economic & 0.959 & 0.791 & +0.168 \\
Social & 0.686 & 0.577 & +0.109 \\
\bottomrule
\end{tabular}
\captionof{table}{B2 confidence on correct and incorrect MITweet predictions.}
\label{tab:confidence}
\end{center}

Table~\ref{tab:confidence} shows that correct B2 predictions are more confident than errors for every projected MITweet label. This is useful because an abstaining classifier only helps when confidence is meaningfully different on correct and incorrect cases. The largest gap appears for economic direction, which is also the weakest RoBERTa result.

\begin{center}
    \centering
    \includegraphics[width=\columnwidth]{figure03_robustness_panel.pdf}
    \captionof{figure}{Robustness diagnostics under noise and counterfactual edits. The figure checks whether the evaluation code records degradation patterns, not whether B5 or B6 has been fully validated on Reddit-scale data.}
    \label{fig:robustness}
\end{center}

The B5 and B6 comparison was verified on a small synthetic smoke setup. The smoke run used three seeds and 20 shared validation posts for the paired bootstrap comparison. It is useful for checking that prediction export, supported-label macro-F1, paired bootstrap testing, and error analysis run end-to-end. It is not a final claim about model quality. Table~\ref{tab:smoke} shows that B6 underperforms B5 under this constrained setup. Relevance is similar, but B6 is weaker on the ideology heads.

\begin{center}
\centering
\small
\begin{tabular}{lcc}
\toprule
Metric & B5 & B6 \\
\midrule
Macro-F1 mean & 0.356 & 0.154 \\
Relevance & 0.467 & 0.473 \\
Economic & 0.322 & 0.194 \\
Social & 0.222 & 0.190 \\
\bottomrule
\end{tabular}
\captionof{table}{Three-seed smoke-run results for the flat and hierarchical transformer variants.}
\label{tab:smoke}
\end{center}

\begin{center}
    \centering
    \includegraphics[width=\columnwidth]{figure02_ablation_deltas.pdf}
    \captionof{figure}{Ablation delta chart for model components and diagnostics. Positive bars indicate settings that improved the supported metric in the current artifact set.}
    \label{fig:ablation}
\end{center}

\begin{center}
    \centering
    \includegraphics[width=\columnwidth]{figure04_error_profile_panel.pdf}
    \captionof{figure}{Error profile separating failure types and content tags. The panel is meant to show where the system fails, rather than only reporting one aggregate score.}
    \label{fig:error-profile}
\end{center}
\FloatBarrier

\section{Limitations}

The full Reddit weak-supervision experiment could not be completed because Politosphere access was blocked in the available environment. As a result, B3 through B6 should be treated as implemented and smoke-tested rather than fully evaluated on the intended Reddit scale.

The MITweet experiment is real-data evidence, but it only covers the labels that can be cleanly projected from MITweet: relevance, economic direction, and social direction. It does not fully evaluate target, stance, frame, intensity, or ambiguity in the same way as the Reddit annotation schema.

The visualizations summarize current artifacts, but some robustness and counterfactual consistency values are diagnostics rather than direct B5 and B6 inference on every edited example. The next experimental step is class-weighted or focal loss for the ideology heads, followed by a full real-data rerun once the Reddit weak-supervision data is reachable.

\section{Conclusion}

This paper ends with a mixed but useful result. The real MITweet experiment supports the basic claim that ideology prediction from social media text is possible, but it does not support a simple claim that RoBERTa is always better than a lighter baseline. RoBERTa improves social-direction prediction, where the labels are less dominated by one class. It performs worse on economic direction because the projected MITweet labels are mostly neutral, and the model learns that majority pattern too strongly. The balanced TF-IDF baseline is weaker as a representation model, but its class weighting makes it more reliable on the minority economic classes.

The B5 and B6 branch is best understood as an implementation milestone. It shows that the code can run ingestion, weak-label construction, model training, prediction export, paired bootstrap testing, calibration checks, and error analysis. It does not yet show that the hierarchical ideology representation improves final accuracy. That distinction matters because otherwise the project would overstate the smoke-run figures. The current evidence says that the evaluation framework is ready, but the full weak-supervision claim still depends on rerunning the system with accessible Reddit data and imbalance-aware training.

The main lesson is that multidimensional ideology modeling needs both better labels and better evaluation discipline. A single aggregate score hides the central failure case in this experiment: a model can appear strong while missing minority ideological directions. Future work should keep the model ladder, but judge it with class-weighted or focal losses, soft-label training for annotator disagreement, and topic or community shift tests. Those changes would make the experiment match the original goal more closely: predicting ideology as a structured signal instead of treating political text as one flat label.

\clearpage
\begin{thebibliography}{7}

\bibitem[Liu et~al.(2023)]{liu-etal-2023-ideology}
Songtao Liu, Ziling Luo, Minghua Xu, Lixiao Wei, Ziyao Wei, Han Yu, Wei Xiang, and Bang Wang.
\newblock Ideology takes multiple looks: A high-quality dataset for multifaceted ideology detection.
\newblock In \emph{Proceedings of the 2023 Conference on Empirical Methods in Natural Language Processing}, pages 4200--4213, 2023.

\bibitem[Chen et~al.(2023)]{chen-etal-2023-ideology}
Chen Chen, Dylan Walker, and Venkatesh Saligrama.
\newblock Ideology prediction from scarce and biased supervision: Learn to disregard the ``what'' and focus on the ``how''!
\newblock In \emph{Proceedings of the 61st Annual Meeting of the Association for Computational Linguistics}, pages 9529--9549, 2023.

\bibitem[Hofmann et~al.(2022)]{hofmann-etal-2022-politosphere}
Valentin Hofmann, Hinrich Sch{"u}tze, and Janet~B. Pierrehumbert.
\newblock The Reddit Politosphere: A large-scale text and network resource of online political discourse.
\newblock In \emph{Proceedings of the International AAAI Conference on Web and Social Media}, volume 16, pages 1259--1267, 2022.

\bibitem[Colacrai et~al.(2024)]{colacrai-etal-2024-politicalcompass}
Ernesto Colacrai, Federico Cinus, Gianmarco De~Francisci Morales, and Michele Starnini.
\newblock Navigating multidimensional ideologies with Reddit's Political Compass: Economic conflict and social affinity.
\newblock \emph{EPJ Data Science}, 2024.

\bibitem[Liu et~al.(2019)]{liu2019roberta}
Yinhan Liu, Myle Ott, Naman Goyal, Jingfei Du, Mandar Joshi, Danqi Chen, Omer Levy, Mike Lewis, Luke Zettlemoyer, and Veselin Stoyanov.
\newblock RoBERTa: A robustly optimized BERT pretraining approach.
\newblock arXiv:1907.11692, 2019.

\bibitem[Chang and Masterson(2020)]{chang-masterson-2020-wordorder}
Charles Chang and Michael Masterson.
\newblock Using word order in political text classification with long short-term memory models.
\newblock \emph{Political Analysis}, 28(3):395--411, 2020.

\bibitem[Leonardelli et~al.(2025)]{leonardelli-etal-2025-lewidi}
Elisa Leonardelli, Silvia Casola, Siyao Peng, Giulia Rizzi, Valerio Basile, Elisabetta Fersini, Diego Frassinelli, Hyewon Jang, Maja Pavlovic, Barbara Plank, and Massimo Poesio.
\newblock LeWiDi-2025 at NLPerspectives: Third edition of the learning with disagreements shared task.
\newblock In \emph{Proceedings of the 4th Workshop on Perspectivist Approaches to NLP}, pages 182--195, 2025.

\end{thebibliography}

\end{document}
